% 復習3　先行研究の探し方
\session{3}{1}{90分}{先行研究の探し方}{AI検索の活用}{}

\goals{
  \item AIを使って関連分野と検索キーワード（日本語と英語）を広げられる
  \item 学術データベースや公的統計で文献を探し、実在と内容を確認できる
  \item 文献リストに必要な情報（著者、発行年、題名、掲載誌または出版社、巻号と頁、DOIまたはURL）を揃えられる
}

\fillin{自分のリサーチ・クエスチョン}

\work[60mm]{ワーク1}{探す場所の見当をつける}
\begin{promptbox}[探す場所の見当をつける]
次のリサーチ・クエスチョンに関連する学術分野と、文献検索に使うキーワードを日本語と英語でそれぞれ10個ずつ挙げてください。文献名や著者名は挙げないでください。問い：（あなたの問い）
\end{promptbox}
\fillin{関連する学術分野}
\begin{wstabh}{colspec={X[l,m] X[l,m] X[l,m] X[l,m]},row{2-Z}={ht=8mm}}
  日本語キーワード & 英語キーワード & 日本語キーワード & 英語キーワード \\
  & & & \\ & & & \\ & & & \\ & & & \\ & & & \\
\end{wstabh}

\work[50mm]{ワーク2}{データベースで検索する}
\instr{CiNii Research、Google Scholar、J-STAGE、図書館の蔵書検索、官公庁の統計・白書などを使う。検索日と件数も記録する。}
\begin{wstabh}{colspec={Q[l,wd=28mm,m] X[2,l,m] Q[c,wd=14mm,m] X[3,l,m]},row{2-Z}={ht=10mm}}
  データベース & 検索語 & 件数 & 関係の深そうな文献 \\
  & & & \\ & & & \\ & & & \\ & & & \\ & & & \\
\end{wstabh}

\work[60mm]{ワーク3}{AIが挙げた文献の実在を確かめる}
\instr{AIに「文献を挙げて」と頼み、挙がったものをデータベースで確認する。}
\begin{wstabh}{colspec={X[3,l,m] Q[c,wd=13mm,m] Q[c,wd=13mm,m] X[2,l,m]},row{2-Z}={ht=10mm}}
  AIが挙げた文献 & 実在\newline ○× & 内容\newline 一致 & 確認先・違っていた点 \\
  & & & \\ & & & \\ & & & \\ & & & \\ & & & \\
\end{wstabh}
{\small 集計：挙がった文献 \underline{\hspace{10mm}} 件のうち、実在しない \underline{\hspace{10mm}} 件、内容が違う \underline{\hspace{10mm}} 件}\par

\work[80mm]{まとめ}{文献リスト（実在を確認した5件）}
\begin{wstabh}{colspec={Q[c,wd=6mm,m] X[2,l,m] Q[c,wd=10mm,m] X[3,l,m] X[3,l,m] X[2,l,m]},row{2-Z}={ht=15mm}}
  & 著者 & 年 & 題名 & 掲載誌・出版社、巻号・頁 & DOI・URL \\
  1 & & & & & \\ 2 & & & & & \\ 3 & & & & & \\ 4 & & & & & \\ 5 & & & & & \\
\end{wstabh}

\begin{nexttime}
  \item 実在を確認した文献5件のリストを完成させる。うち1件は本文を入手する（復習4で読む）
  \item AIが挙げた文献のうち、実在しなかった・内容が違ったものの数と例を記録しておく（上のワーク3）
\end{nexttime}
\aimemo
